\subsection{Общая структура программы}
\label{subsec:program-structure}

Структура проектируемой программы разделяет описание нагрузки, полезную обработку данных,
  управление исполнением и получение результата. Такое устройство реализует требования
  раздела~\ref{sec:experiment-design}: один тип задачи можно запускать в разных режимах,
  сохраняя его входные данные и алгоритм. Взаимодействие компонентов определяется передаваемыми
  конфигурациями, задачами и результатами, а не особенностями конкретного режима многопоточности.

Программа представляет собой один исполняемый файл. Внутри него выделяются модули конфигурации,
  подготовки сценария, задач, исполнения, управления экспериментом и вывода.
  Предлагаемое распределение по пакетам приведено в таблице~\ref{tab:program-modules}.
  Каждый пакет отвечает за ограниченную часть обработки; самостоятельные процессы для этих
  компонентов не требуются. Пакеты конфигурации, сценариев, задач, исполнения, эксперимента
  и вывода размещаются внутри каталога \texttt{internal}; в таблице их имена даны относительно него.

\begin{table}[htbp]
  \centering
  \caption{Компоненты программы и их назначение}
  \label{tab:program-modules}
  \begin{tabularx}{\textwidth}{|p{4.5cm}|X|}
    \hline
    Компонент & Назначение \\
    \hline
    \texttt{cmd/experiment} & Точка входа, выбор запуска и завершение приложения \\
    \hline
    \texttt{config} & Представление и проверка параметров \\
    \hline
    \texttt{scenario} & Подготовка данных и состава нагрузки \\
    \hline
    \texttt{task} & Общий контракт задачи и результата \\
    \hline
    \texttt{tasks/logs} & Разбор записей и расчёт статистики \\
    \hline
    \texttt{tasks/compress} & Сжатие независимых блоков \\
    \hline
    \texttt{executor} & Выполнение задач в выбранном режиме \\
    \hline
    \texttt{experiment} & Управление запуском, измерением и проверкой \\
    \hline
    \texttt{output} & Формирование итоговой записи о запуске \\
    \hline
  \end{tabularx}
\end{table}

Модуль конфигурации хранит типизированное описание запуска. Общие параметры задают режим,
  число горутин, \texttt{GOMAXPROCS}, ёмкость очереди и объём данных. Параметры конкретного
  типа задают размер пакета, правила фильтрации, уровень сжатия и другие свойства полезной работы.
  Проверка конфигурации выполняется до подготовки основного набора: отрицательный размер,
  неизвестный режим или несовместимое сочетание настроек приводят к отклонению запуска.
  После проверки используется согласованный набор значений, включая явно определённые значения
  по умолчанию.

Модуль сценария преобразует конфигурацию в подготовленную нагрузку. Он создаёт исходные данные,
  определяет границы блоков, формирует точное число задач каждого класса и задаёт порядок выдачи.
  Описание элемента содержит идентификатор, класс, конфигурацию типа и ссылку на входной участок.
  Это позволяет сохранять общий неизменяемый набор данных вместо копирования всего входа
  для каждого элемента. Последний блок может быть меньше остальных, но также входит в обработку.

Данные и список задач создаются до измерения. Раздельные \texttt{data\_seed}
  и \texttt{order\_seed} управляют содержимым и перестановкой списка, как определено
  в подразделе~\ref{subsec:reproducible-workloads}. Для следующего повторения создаётся
  новое изменяемое состояние. Подготовленное описание нагрузки позволяет воспроизвести состав
  задач независимо от того, как планировщик распределит их между горутинами.

Общий контракт задачи задаётся интерфейсом \texttt{Task} с методом
  \texttt{Run(ctx context.Context) (Result, error)}. Конкретная структура хранит конфигурацию
  и входные данные, полученные конструктором. Вызов \texttt{Run} выполняет одну единицу полезной
  работы и возвращает её результат. Задача не выбирает модель потоков, не изменяет
  \texttt{GOMAXPROCS} и не управляет всей очередью. Благодаря этому различие режимов сосредоточено
  в исполнителе, а полезные алгоритмы остаются общими.

Модуль обработки журналов создаёт \texttt{LogsTask}. Его вход представляет пакет записей,
  выход — частичную статистику либо обновление явно переданного общего накопителя.
  При локальной политике итоговая статистика получается объединением частичных результатов задач.
  При общей политике доступ к накопителю защищается одним ресурсом синхронизации сценария,
  а итоговая статистика считывается после завершения всех обновлений. В этом случае результаты
  задач подтверждают обработанный объём, но повторно не прибавляются к общей статистике.
  Для сопоставления режимов сохраняется одинаковая частота обновлений: блокировка на каждую запись
  и блокировка на целый пакет задают разные нагрузки.

Модуль сжатия создаёт \texttt{CompressTask}, работающую с независимым блоком байтов.
  Конфигурация фиксирует уровень сжатия и свойства блока. Изменяемое состояние преобразования
  принадлежит отдельной задаче; результат содержит сжатый блок и обработанный объём.
  Конструктор подготавливает описание задачи, а рабочие операции и выделения памяти внутри
  \texttt{Run} входят в измеряемую обработку. Восстановление исходного блока относится
  к последующей проверке корректности.

Исполнитель получает готовый список задач и предоставляет три режима. В \texttt{seq}
  элементы обрабатываются последовательно. В \texttt{mn} создаётся заданное число горутин,
  получающих задачи через ограниченную очередь. В \texttt{pinned} используется тот же пул
  с закреплением каждой горутины обработки за потоком ОС. Закрепление выполняется до сообщения
  о готовности, чтобы его подготовительная стоимость не попадала во время обработки.
  Режимы \texttt{mn} и \texttt{pinned} используют одинаковые правила выдачи и приёма результатов.

Горутина обработки последовательно получает задачи, вызывает \texttt{Run} и передаёт результат
  вместе с идентификатором. Завершение всех задач определяется по факту их обработки,
  а не только по опустошению очереди: после выдачи последнего элемента часть горутин ещё может
  выполнять работу. На момент остановки измерения все полезные результаты должны быть доступны
  и объединены. Освобождение закреплений и окончательное завершение служебных горутин выполняются
  после этого и не подменяют границу измерения.

Модуль управления экспериментом связывает остальные компоненты. Он инициирует подготовку,
  ожидает готовности исполнителей, задаёт общий старт и фиксирует окончание обработки.
  Этот модуль также организует объединение результатов, запускает проверку и формирует описание
  завершившегося эксперимента. Измерение полного сценария находится в одном месте, поэтому
  изменение типа задачи не меняет смысл измеряемого времени.

Результаты связываются с уникальными идентификаторами задач, а не с порядком завершения.
  При приёме проверяются принадлежность идентификатора исходному списку и отсутствие ранее
  принятого результата с тем же идентификатором. После обработки множество полученных
  идентификаторов должно совпадать с множеством задач сценария. Для сжатия дополнительно
  сохраняется соответствие результата исходному блоку, а статистика объединяется по заданным
  правилам. Исполнитель обеспечивает однократный вызов \texttt{Run} для каждой задачи;
  проверка идентификаторов обнаруживает пропущенные и повторно полученные результаты.

Модуль вывода получает уже сформированную запись о запуске и сохраняет параметры, длительность,
  объём полезной работы и статус проверки. Запись выполняется после остановки измерения.
  Подробный вывод из цикла обработки не используется в основной серии, поскольку он добавлял бы
  операции синхронизации и ввода-вывода. Формирование итоговой записи не должно зависеть от типа
  исполнителя: режим является одним из её полей.

При ошибке задачи управление экспериментом прекращает выдачу новых элементов и передаёт сигнал
  отмены. Исполнитель завершает участвующие горутины и снимает закрепления, а модуль вывода
  сохраняет неуспешный статус с указанием причины. Частичный результат не считается завершённой
  обработкой. Для нового типа задачи требуется определить конфигурацию, конструктор,
  \texttt{Run} и способ проверки; механизм запуска и учёта времени сохраняется.

% TODO: при появлении кода сверить имена пакетов и контракты с реализацией.
% Конфигурация MVP задаётся структурами Go. YAML подключить поверх того же внутреннего представления.
% Метрики, трейсы, Grafana и pprof — отдельная итерация; не описывать их как реализованные компоненты.

\subsection{Функциональная схема и алгоритм программы}
\label{subsec:program-schemes}

% TODO: добавить функциональную схему: конфигурация -> подготовка сценария -> исполнитель -> задачи
% -> объединение результатов -> проверка -> итоговый лог. Отдельно показать управление измерением.
% TODO: добавить блок-схему запуска с ветвями seq/mn/pinned, готовностью горутин, стартом и остановкой.
% Отразить ошибки, отмену, снятие закреплений и разделение подготовки, обработки и проверки.
% Схемы сопровождаются подписями и ссылками из текста после их добавления.

\subsection{Условия проведения эксперимента}
\label{subsec:experiment-conditions}

Для исследования выбраны две платформы, характеристики которых приведены
  в подразделах~\ref{subsec:test-systems} и~\ref{subsec:software-platforms}.
  Система x86 использует Intel Core i5-13500H и Ubuntu Linux, система ARM — MacBook Pro
  16 дюймов с Apple M1 Pro, десятью ядрами CPU и macOS 26.
  Для каждой системы предусматривается нативный исполняемый файл: x86-64 для Intel
  и arm64 для Apple. Это исключает добавление трансляции команд другой архитектуры
  в основную пару сравнений.

Основная единица сравнения — два запуска одного сценария на одной машине в режимах
  \texttt{mn} и \texttt{pinned}. Внутри пары совпадают версия исходного кода, инструменты сборки,
  входные данные, список задач, параметры полезной обработки, число горутин, \texttt{GOMAXPROCS}
  и ёмкость очереди. Последовательный запуск выполняет тот же объём операций и позволяет
  отдельно оценить эффект параллелизации.

Состав программной среды связывается с результатом каждого запуска. В него входят версия ОС,
  версия ядра Linux для Ubuntu, версия Go, архитектура исполняемого файла и параметры сборки.
  Число доступных логических CPU фиксируется отдельно от заданного \texttt{GOMAXPROCS}.
  Это различает свойства стенда и ограничение, установленное программой для конкретного опыта.

% TODO: дополнить описание и итоговую таблицу фактическими версиями Ubuntu, ядра, macOS и Go.
% TODO: добавить модель x86-ноутбука, объём и конфигурацию RAM, параметры сборки и идентификатор кода.
% Эти данные вводятся после проверки стендов; неизвестные значения не выводить в текст заглушками.

Условия питания и фоновой нагрузки сохраняются сопоставимыми внутри серии. Питание от сети,
  режим энергопотребления и состояние системы входят в описание стенда. Нагрузка сторонних
  приложений ограничивается, а подготовка между запусками не должна создавать дополнительную
  вычислительную работу одновременно с измерением. Для ноутбуков учитывается возможность
  изменения частоты из-за нагрева, поэтому порядок запусков и промежутки между ними являются
  частью условий серии.

% TODO: зафиксировать фактические настройки питания и порядок стабилизации температуры на каждом стенде.

Ни один из режимов не задаёт привязку к конкретному ядру. Размещение системных потоков
  по производительным и энергоэффективным ядрам остаётся под управлением ОС.
  \texttt{GOMAXPROCS} задаётся явно, но не определяет, какой класс ядер будет использоваться.
  Сопоставление результатов поэтому учитывает неоднородность обеих систем и не приравнивает
  одинаковое число горутин к одинаковой аппаратной производительности.

Граница времени обработки начинается с общего сигнала старта после подготовки данных и горутин.
  В интервал входят выдача задач, вызовы \texttt{Run}, ожидания, предусмотренные нагрузкой,
  и объединение результатов. Получение последнего объединённого результата завершает интервал.
  Проверка полезного выхода, запись итогового лога и завершение инфраструктуры выполняются
  после остановки. Для всех режимов сохраняется один смысл начала и окончания измерения.

Длительность обработки измеряется как прошедшее время по монотонным часам средствами
  \texttt{time.Now} и \texttt{time.Since} с сохранением монотонной составляющей отметки
  старта~\cite{go_time_package}. Этот показатель включает ожидания и отличается от процессорного
  времени, затраченного программой.

При необходимости предварительного прогрева он выполняется на отдельном изменяемом состоянии.
  Его результаты не входят в измеряемую обработку. Для независимого повторения запускается
  новый процесс, чтобы накопленные системные потоки и состояние очередей предыдущего запуска
  не переносились непосредственно в следующий. Это не устраняет влияние кэшей, температуры
  и фоновой активности ОС, которое учитывается повторными измерениями.

До основной серии проверяется корректность обработки: статистика журналов сопоставляется
  с последовательным эталоном, а сжатые блоки восстанавливаются и сравниваются с исходными.
  Отдельно проверяется отсутствие гонок. Инструментированная сборка для такой проверки
  не используется для измерения производительности.

Режим закрепления проверяется отдельно от основного измерения времени. Число живых потоков ОС
  сопоставляется с числом закреплённых горутин с учётом служебных потоков среды выполнения.
  Проверка привязки горутины к потоку и наблюдение числа потоков решают разные задачи:
  одно только увеличение их количества не подтверждает правильность всех действий исполнителя.
  Средства диагностики и момент наблюдения фиксируются вместе с условиями проверки.

% TODO: после реализации описать фактическую проверку закрепления и способы наблюдения потоков
% на Linux и macOS. Не заменять число потоков значением runtime.NumGoroutine.

Профилирование и подробные трассировки относятся к отдельным диагностическим запускам.
  В основной серии сохраняется одинаковый набор измерительных действий в обоих режимах.
  Это позволяет оценивать длительность общей обработки без различающейся нагрузки от наблюдения.
  Показатели памяти и потоков сопровождаются способом сбора, если они используются в сравнении.

Повторение опыта на второй платформе проверяет сохранение тенденций. Поскольку одновременно
  меняются процессор, память, ОС и условия охлаждения, межплатформенная разница относится
  к двум системам целиком. Причины различий между моделями исследуются прежде всего внутри
  каждой системы при сохранённых условиях запуска.

\subsection{Серии измерений и порядок проведения исследования}
\label{subsec:measurement-series}

% TODO: описать пилот, окончательную матрицу P/W/размеров блоков и условия каждой серии.
% Включить однотипные задачи, общую/локальную статистику, смеси классов и разные order_seed.
% Ожидание рассматривать отдельной серией. Указать объём работы, повторения и чередование режимов.
% Согласовать состав серий с фактически реализованными типами задач и параметрами.

\subsection{Результаты измерений и их анализ}
\label{subsec:experiment-results}

% TODO: привести исходные и сводные результаты после измерений; без вымышленных значений.
% Сначала сравнить seq/mn/pinned внутри каждого стенда, затем сопоставить тенденции двух платформ.
% Добавить таблицы и графики времени, пропускной способности, ускорения и отношения Tpinned/Tmn.
% При наличии данных разобрать память, потоки ОС и задержки по классам задач с учётом разброса.
% Проверить гипотезы раздела 3, объяснить наблюдения и указать ограничения выводов.
